\chapter{Gravitation and orbits}\label{ch:grav}

Newton's inverse-square law is the first ``field'' in the book: a
$1/r^{2}$ fall-off that runs from falling apples to GPS. This chapter
does the useful part --- the shell theorem, escape and orbital speeds,
and orbits as bookkeeping on $E$ and $L$.

\section{Newton's law and the field picture}\label{sec:grav-law}

\begin{equation}
  F=G\frac{m_1m_2}{r^{2}},\qquad
  \mathbf g(\mathbf r)=-G\frac{M}{r^{2}}\hat r,\label{eq:grav}
\end{equation}
where $G=6.674\times10^{-11}$ SI. The field picture separates cause and
test: $M$ builds $\mathbf g$, and any $m$ placed in it feels
$m\mathbf g$. Cavendish weighed the Earth with this: $g=GM/R^{2}$.

\section{The shell theorem}\label{sec:grav-shell}

Newton invented the shell theorem to justify treating the Earth as a
point: outside a spherical shell its field is that of all mass at the
centre; inside it the field is \emph{zero} (figure~\ref{fig:shell}).
Consequence: inside a uniform planet only the mass beneath you pulls
(figure~\ref{fig:gr}).

\begin{figure}[htbp]\centering
\begin{tikzpicture}[scale=0.9]
\draw[very thick] (0,0) circle (1.8);
\draw[fill=accent!15] (0,0) circle (1.8);
\fill[white] (0,0) circle (0.8);
\draw[very thick] (0,0) circle (0.8);
\fill (0.3,1.3) circle (0.04) node[right,font=\scriptsize]{$m$ inside};
\fill (-1.0,0.3) circle (0.04) node[right,font=\scriptsize]{$m$ in shell};
\node[font=\scriptsize,align=left] at (2.4,0.8) {shell theorem:\\inside $\to$ mass inside only\\in-shell $\to$ zero};
\end{tikzpicture}

\caption{Shell theorem: mass above the point contributes zero; mass below
acts from the centre. Gravitational depth-charge accounting.}
\label{fig:shell}
\end{figure}

\begin{figure}[htbp]\centering
\begin{tikzpicture}
\begin{axis}[width=8.5cm,height=5.2cm,xlabel={$r/R_\oplus$},ylabel={$g(r)$ (rel.)},
xmin=0,xmax=3,ymin=0,ymax=1.3,font=\small]
\addplot[very thick,accent,domain=0:1] {x};
\addplot[very thick,accent,domain=1:3] {1/x^2};
\addplot[mark=*] coordinates {(1,1)} node[above right,font=\scriptsize]{surface};
\node[font=\scriptsize] at (axis cs:0.5,0.35) {uniform: $g\propto r$};
\node[font=\scriptsize] at (axis cs:2.1,0.5) {$g\propto 1/r^{2}$};
\end{axis}
\end{tikzpicture}

\caption{$g(r)$ inside a uniform Earth (linear) and outside
($1/r^{2}$) --- the two regimes of the shell theorem.}
\label{fig:gr}
\end{figure}

\begin{example}[$g(r)$ inside the Earth]\label{ex:ginside}
Find $g$ at radius $r<R$ in a uniform Earth.
\par\smallskip
\textbf{Strategy.} Only mass below: $M(r)=M(r/R)^{3}$.
\begin{solution}
$g(r)=GM(r)/r^{2}=GM r/R^{3}=g_{0}\,r/R$: gravity drops linearly to zero
at the centre. The fastest ``tunnel'' is down to $g=0$ at $r=0$, not up.
\checkmark
\end{solution}
\end{example}

\section{Gravitational potential energy}\label{sec:grav-u}

The work integral over \eqref{eq:grav} gives
\begin{equation}
  U(r)=-G\frac{Mm}{r},\qquad \Delta U=GMm\!\left(\tfrac1{r_1}-\tfrac1{r_2}\right).
  \label{eq:ug}
\end{equation}
The zero at infinity is the clean choice: bound systems carry negative
energy. Near the surface it reduces to $mgh$: expand
$\Delta U\approx GMm\,h/R^{2}=mgh$.

\begin{example}[Escape velocity]\label{ex:escape}
Find the speed to leave Earth from the surface and explain why direction
doesn't matter.
\par\smallskip
\textbf{Strategy.} Escape means arriving at infinity with $K\ge0$:
$K_0+U_0\ge0$.
\begin{solution}
$\tfrac12mv_{\rm esc}^{2}=GMm/R \Rightarrow
v_{\rm esc}=\sqrt{2GM/R}=11.2$\,km/s for Earth. Energy is scalar ---
aim up, sideways or through a canyon (if it could): the bill is the same.
Compare orbital speed $\sqrt{GM/R}=7.9$\,km/s: escape is
$\sqrt2$ times orbital. \checkmark
\end{solution}
\end{example}

\begin{exercise}[easy]
At what altitude is $v_{\rm esc}$ down to half its surface value? (F=ma)
\end{exercise}

\section{Circular orbits}\label{sec:grav-orbit}

Centripetal balance gives
\begin{equation}
  v_{\rm orb}=\sqrt{GM/r},\quad T=2\pi\sqrt{r^{3}/GM},\quad
  E=-\frac{GMm}{2r},\label{eq:orb}
\end{equation}
and the energy bookkeeping is odd: $K=+GMm/2r$, $U=-GMm/r$, so $E=-K$
--- a satellite that loses energy (drag) \emph{speeds up} as it sinks.

\begin{example}[ISS energy audit]\label{ex:iss}
The ISS ($m=4.2\times10^{5}$\,kg) orbits at $r\approx6800$\,km.
Find $v$, $K$, $U$, $E$ and the energy to raise it $10$\,km.
\par\smallskip
\textbf{Strategy.} Equation~\eqref{eq:orb} line by line; the last answer
is $\Delta E>0$ for a \emph{larger} $r$.
\begin{solution}
$v=\sqrt{GM/r}=\sqrt{3.99\times10^{14}/6.8\times10^{6}}=7.66$\,km/s.
$K=\tfrac12mv^{2}=1.23\times10^{13}$\,J, $U=-2.47\times10^{13}$\,J,
$E=-1.23\times10^{13}$\,J.
Raising $r$ by $10$\,km: $\Delta E=GMm\,\Delta r/2r^{2}
\approx9.1\times10^{9}$\,J --- nine GJ buys ten kilometres.
\checkmark
\end{solution}
\end{example}

\section{Kepler's laws, derived}\label{sec:grav-kepler}

\begin{enumerate}
\item \emph{Ellipses}: the $1/r^{2}$ law's bound orbits are conics ---
stated here and used everywhere below.
\item \emph{Equal areas}: $dA/dt=L/2m$ is constant --- conservation of
angular momentum, nothing else (figure~\ref{fig:kepler}).
\item \emph{$T^{2}\propto a^{3}$}: from \eqref{eq:orb},
$T^{2}=4\pi^{2}a^{3}/GM$ --- our code check recovers it to a constant.
\end{enumerate}

\begin{example}[Perihelion vs aphelion speed]\label{ex:kepler2}
An orbit has $e=0.6$: find $v_{\rm peri}/v_{\rm apo}$.
\par\smallskip
\textbf{Strategy.} Kepler II is $L$-conservation: $rv$ equal at the two
axis ends.
\begin{solution}
$r_p=a(1-e)=0.4a$, $r_a=a(1+e)=1.6a$;
$v_p/v_a=r_a/r_p=(1+e)/(1-e)=4$. Four-to-one speed ratio --- comets'
sun-grazing sprints are this exact arithmetic. \checkmark
\end{solution}
\end{example}

\begin{figure}[htbp]\centering
\begin{tikzpicture}[scale=1.3]
\draw[very thick,accent] (0,0) ellipse (2.4 and 1.3);
\fill[amber] (-1.8,0) circle (0.05) node[below,font=\scriptsize]{sun};
\fill (-1.8,0) circle (0.001);
\draw[thick] (-1.8,0) -- (2.4,0);
\draw[thick] (-1.8,0) -- (1.2,1.13);
\draw[thick] (-1.8,0) -- (-0.8,-1.18);
\draw[thick] (-1.8,0) -- (-2.3,-0.36);
\draw[fill=accent!15,draw=none] (-1.8,0) -- (2.4,0) arc (0:43:2.4 and 1.3) -- cycle;
\draw[fill=accent!15,draw=none] (-1.8,0) -- (-0.8,-1.18) arc (200:235:2.4 and 1.3) -- cycle;
\node[font=\scriptsize] at (0.5,-0.5) {equal areas, equal times};
\node[font=\scriptsize] at (-1.0,-1.6) {perihelion fast, aphelion slow};
\end{tikzpicture}

\caption{Kepler's second law is $L$-conservation drawn: equal areas in
equal times.}
\label{fig:kepler}
\end{figure}

\begin{example}[Hohmann transfer bookkeeping]\label{ex:hohmann}
Move a satellite from $r_1=R$ to $r_2=2R$ (Earth radii). Find $\Delta v$
at each burn: $v_1=\sqrt{GM/R}=7.9$\,km/s, $v_2=\sqrt{GM/2R}=5.6$\,km/s.
\par\smallskip
\textbf{Strategy.} Transfer ellipse has $a_t=(r_1+r_2)/2=1.5R$; its
peri/apo speeds come from energy.
\begin{solution}
Vis-viva $v^{2}=GM(2/r-1/a_t)$: at $r_1$,
$v_{t1}=\sqrt{GM(2/R-2/3R)}=\sqrt{4GM/3R}=9.1$\,km/s
$\Rightarrow\Delta v_1=+1.2$\,km/s; at $r_2$,
$v_{t2}=\sqrt{GM(1/R-2/3R)}=4.6$\,km/s
$\Rightarrow\Delta v_2=+1.0$\,km/s. Total $2.2$\,km/s ---
figure~\ref{fig:hohmann}. \checkmark
\end{solution}
\end{example}

\begin{figure}[htbp]\centering
\begin{tikzpicture}[scale=1.1]
\draw[very thick] (0,0) circle (0.8);
\draw[very thick] (0,0) circle (2.0);
\fill[amber] (0,0) circle (0.07) node[right=2pt,font=\scriptsize]{};
\draw[very thick,dashed,accent] (0.8,0) arc (0:180:1.4 and 1.0);
\draw[->,thick,accent] (0.8,0) -- (1.15,0.35) node[right,font=\scriptsize]{$\Delta v_1$};
\draw[->,thick,accent] (-2.0,0) -- (-1.55,-0.35) node[left,font=\scriptsize]{$\Delta v_2$};
\node[font=\scriptsize] at (0,0.95) {$r_1$};
\node[font=\scriptsize] at (0,-2.2) {$r_2$};
\node[font=\scriptsize,align=left] at (-1.3,1.2) {transfer ellipse:\\$a_t=(r_1+r_2)/2$};

\end{tikzpicture}

\caption{Hohmann transfer: two burns, one ellipse, the cheapest lane
between circular orbits.}
\label{fig:hohmann}
\end{figure}

\begin{figure}[htbp]\centering
\begin{tikzpicture}
\begin{axis}[width=8cm,height=4.8cm,xlabel={$a$ (AU)},ylabel={$E$ (rel.)},
xmin=0.4,xmax=3,ymin=-1.3,ymax=0.3,font=\small]
\addplot[very thick,accent,domain=0.45:3,samples=60] {-0.5/x};
\addplot[dashed,amber] coordinates {(0.4,-1.25) (3,-1.25)} node[below,pos=0.5,font=\scriptsize]{inner orbit};
\addplot[dashed,amber] coordinates {(0.4,-0.25) (3,-0.25)} node[above,pos=0.5,font=\scriptsize]{outer orbit};
\node[font=\scriptsize,align=left] at (axis cs:1.7,-0.7) {transfer pays\\the gap $\Delta E$};
\end{axis}
\end{tikzpicture}

\caption{Energy vs semi-major axis: $E=-GMm/2a$, so raising an orbit
means paying the gap.}
\label{fig:ea}
\end{figure}

\section{$E$ and $L$: the orbit's conserved labels}\label{sec:grav-cons}

An orbit is fully labelled by $(E,L)$; the effective potential
$U_{\rm eff}=L^{2}/2mr^{2}-GMm/r$ (figure~\ref{fig:ur}) shows why:
$L^{2}/2mr^{2}$ pushes you out at short range, $-GMm/r$ pulls in at
long, and the bound orbit oscillates in the bowl --- perihelion and
aphelion are the turning points of chapter~7's $U(x)$ grammar.

\begin{figure}[htbp]\centering
\begin{tikzpicture}
\begin{axis}[width=8.5cm,height=5.2cm,xlabel={$r$},ylabel={energy},
xmin=0.4,xmax=4,ymin=-2.2,ymax=1.2,font=\small,legend pos=south east]
\addplot[very thick,accent,domain=0.45:4,samples=80] {-1/x};
\addlegendentry{$U=-GMm/r$}
\addplot[very thick,amber,domain=0.45:4,samples=80] {0.5/x^2 - 1/x};
\addlegendentry{$U_{\rm eff}$ ($L>0$)}
\addplot[dashed] coordinates {(0.4,0) (4,0)};
\node[font=\scriptsize] at (axis cs:1.0,-0.55) {bound};
\node[font=\scriptsize] at (axis cs:3.0,0.35) {free};
\end{axis}
\end{tikzpicture}

\caption{$U(r)$ and $U_{\rm eff}(r)$ for nonzero $L$: the centrifugal
barrier makes bound states possible.}
\label{fig:ur}
\end{figure}

\begin{lstlisting}[language=Python,caption={\texttt{code/ch08\_orbit\_verlet.py} --- velocity-Verlet orbit; $E,L$ drift and Kepler check.},label={lst:orbit}]
        ax_, ay_ = -GM*x/r3, -GM*y/r3
        x += vx*dt + 0.5*ax_*dt*dt
        y += vy*dt + 0.5*ay_*dt*dt
        ax2, ay2 = -GM*x/(r_new**3), -GM*y/(r_new**3)
        vx += 0.5*(ax_+ax2)*dt
        vy += 0.5*(ay_+ay2)*dt
\end{lstlisting}

Output of a real run (200{,}000 steps each):
\begin{lstlisting}[numbers=none]
a=1.0 e=0.0: drift E=0.000% L=0.000%  T=6.283
a=1.0 e=0.5: drift E=0.000% L=0.000%  T=6.283
a=1.5 e=0.7: drift E=0.000% L=0.000%  T=11.543

Kepler III check (T^2/a^3 should be const):
  a=0.8: T=4.496  T^2/a^3=39.4784
  a=1.0: T=6.283  T^2/a^3=39.4784
  a=1.3: T=9.313  T^2/a^3=39.4784
  a=1.6: T=12.716  T^2/a^3=39.4784
\end{lstlisting}

Velocity-Verlet holds $E$ and $L$ to $<0.001\%$ even at $e=0.7$ ---
the symplectic lesson of chapter~4's energy-check code paying off ---
and recovers $T^{2}=4\pi^{2}a^{3}/GM=39.4784$ to all printed digits.

\section{Satellites, tides and GPS in two paragraphs}\label{sec:grav-sat}

Real orbits pay rent: LEO drag costs tens of m/s per year, geostationary
station-keeping $\sim50$\,m/s/yr, and GPS clocks tick $45$\,\textmu s/day
fast (relativity) that must be designed out. Tides are the
$d/dr$ of gravity: the Moon's pull differs by $\sim7\%$ across Earth's
diameter, stretching oceans toward and away from it --- the third
derivation nobody expected to be useful in a textbook chapter.

\section{In engineering practice: the launch-and-station budget}\label{sec:grav-eng}

Rocket sizing runs this chapter backwards: the mission's $\Delta v$
ledger (launch $\sim9.4$\,km/s with losses, Hohmann legs, plane changes
at $2v\sin(\theta/2)$) divides by the rocket equation from chapter~3 to
mass ratios; orbit choice runs forwards --- LEO cheap, MEO for nav,
GEO parked. Every constellation on Earth is a worked example of
figure~\ref{fig:ea}.

\section{Common mistakes}\label{sec:grav-mistakes}

\begin{tripwire}[$U=mgh$ used past its altitude]
$mgh$ is the tangent-line approximation of \eqref{eq:ug}; at
$h\sim0.1R$ it's already off by 10\%.
\end{tripwire}
\begin{tripwire}[Escape needs a direction]
The $\sqrt{2GM/R}$ bill is in energy --- any direction works except
into the ground.
\end{tripwire}
\begin{tripwire}[``Weightless'' means ``no gravity'']
LEO $g$ is $\sim90\%$ of surface: astronauts float because ship and crew
fall together --- chapter~2's free-fall frame again.
\end{tripwire}
\begin{tripwire}[$E=-K$ sign confusion]
$E<0$ is bound; $-E$ is the energy to escape from that orbit. And a
sinking satellite speeds up --- the sign flips intuition.
\end{tripwire}
\begin{tripwire}[Kepler III constant confusion]
$T^{2}/a^{3}=4\pi^{2}/GM$ depends on the central mass only: same for
all planets of one star, different per star.
\end{tripwire}

\section{Summary}\label{sec:grav-summary}

\begin{table}[ht]\centering
\begin{tabularx}{\linewidth}{@{}lY@{}}
\toprule
\tablehead\thead{Idea} & \thead{One-line version}\\
\midrule
law & $F=GMm/r^{2}$; $\mathbf g=-GM\hat r/r^{2}$ \\
shell & inside uniform planet $g\propto r$; shell interior $=0$ \\
$U$ & $-GMm/r$; $v_{\rm esc}=\sqrt{2GM/R}$ \\
orbit & $v=\sqrt{GM/r}$, $E=-GMm/2r$ \\
Kepler & ellipses; equal areas $=L$; $T^{2}\propto a^{3}$ \\
transfer & Hohmann: ellipse $a_t=(r_1+r_2)/2$, two burns \\
\bottomrule
\end{tabularx}
\caption{Gravitation cheat sheet.}
\label{tab:gravsummary}
\end{table}

\begin{wins}[Definition of done]
You can now: apply the shell theorem, derive escape and orbital speeds,
balance Hohmann transfers, read orbits off $(E,L)$, and run a
conserving orbit integrator.
\end{wins}

\chapternotes{%
The field-first order follows MIT 8.01SC \cite{ocw801} and AP Physics C
\cite{apced}; the Hohmann numbers are standard mission-analysis figures
(cf.\ Morin \cite{morin} and Feynman vol.~I \cite{feynman1} for the
$C_3$/transfer bookkeeping). The integrator numbers are the listing's
real output.}

\begin{rulebox}[{Watch instead}]
\begin{enumerate}
\item \emph{Kepler's Laws, Elliptical Orbits, Satellites, Orbital Changes}
--- Walter Lewin, MIT 8.01x Lecture 22, 49:09,
\url{https://youtu.be/9wDAm-dlht4} ---
\S\ref{sec:grav-orbit}--\S\ref{sec:grav-cons} at full length (accessed
2026-10-11).
\item \emph{Gravitational potential energy at large distances} --- Khan
Academy, 5:04, \url{https://youtu.be/9Pp_LtKgTQg} ---
\S\ref{sec:grav-u} worked quickly (accessed 2026-10-11).
\end{enumerate}
\end{rulebox}

\section*{Further reading}
\addcontentsline{toc}{section}{Further reading}

\begin{itemize}
\item MIT OpenCourseWare, \emph{8.01SC Classical Mechanics}, gravitation
module --- shell theorem, orbits and Kepler with recitations. At
\url{https://ocw.mit.edu}, course ``8.01SC'' (verified 2026-10-11).
\item OpenStax, \emph{College Physics 2e}, ch.~6 ``Uniform Circular
Motion and Gravitation'' --- free text at \url{https://openstax.org}
(verified 2026-10-11).
\item College Board, \emph{AP Physics C: Mechanics CED}, Unit 4
(gravitation). At \url{https://apstudents.collegeboard.org} (verified
2026-10-11).
\item R.~P. Feynman, \emph{Lectures on Physics}, vol.~I, ch.~7 ``The
Theory of Gravitation'' --- free at Caltech's site.
\item D.~Morin, \emph{Introduction to Classical Mechanics}, ch.~10 ---
orbits and transfers at the harder level; ISBN 978-0521876223.
\item Khan Academy, \emph{Centripetal force and gravitation} --- drills
matching \S\ref{sec:grav-law}--\S\ref{sec:grav-kepler}, at
\url{https://www.khanacademy.org} (verified 2026-10-11).
\end{itemize}

\section*{Additional exercises}
\addcontentsline{toc}{section}{Additional exercises}

\begin{exercise}[easy]
Find the orbital speed and period of a satellite at altitude $2R_E$
above Earth's surface. (AP-C)
\end{exercise}
\begin{exercise}[easy]
Explain why a geostationary orbit's radius is unique, and estimate it
from $T=24$\,h. (AP-C)
\end{exercise}
\begin{exercise}[easy]
Show $v_{\rm esc}=\sqrt2\,v_{\rm orb}$ at the surface and find the
ratio at altitude $R$. (F=ma)
\end{exercise}
\begin{exercise}[easy]
A planet has twice Earth's mass and half its radius: find surface $g$
in Earth-$g$. (F=ma)
\end{exercise}
\begin{exercise}[easy]
Compute the work to lift $1$\,kg from Earth's surface to infinity vs to
orbit radius $2R$. (AP-C)
\end{exercise}
\begin{exercise}[medium]
Estimate the energy per kg for a Hohmann transfer from Earth's orbit to
Mars's ($a=1.5$\,AU) in Solar-System units. (F=ma)
\end{exercise}
\begin{exercise}[medium]
How much does a mine shaft's $g$ drop at depth $R/2$ in a uniform Earth
--- and why is the real drop different? (F=ma)
\end{exercise}
\begin{exercise}[medium]
An elliptical orbit has $a=4R$, $e=0.5$: find peri/apo radii and
speeds, and the period in terms of LEO's. (AP-C)
\end{exercise}
\begin{exercise}[medium]
Explain, with energy signs, why a satellite losing energy to drag moves
faster. (concept)
\end{exercise}
\begin{exercise}[medium]
Find the fraction of Earth's mass ``felt'' by an object at $r=R/2$
inside a uniform Earth, and the resulting $g$. (F=ma)
\end{exercise}
\begin{exercise}[medium]
A plane change of $30^{\circ}$ on a circular LEO orbit: find $\Delta v$
in terms of $v_{\rm orb}$. (engineering)
\end{exercise}
\begin{exercise}[hard]
Show the transfer ellipse of a Hohmann between $r_1$ and $r_2$ has
$T_{\rm tr}=\pi\sqrt{a_t^{3}/GM}$ and compute it for the ISS-to-GEO
case ($r_1=1.05R$, $r_2=6.6R$). (USAPhO)
\end{exercise}
\begin{exercise}[hard]
A satellite in a highly eccentric orbit loses energy only at perigee
drag: argue qualitatively what shrinks and what stays fixed.
(APhO-style)
\end{exercise}
\begin{exercise}[medium]
Two stars of equal mass orbit each other at separation $d$: find the
period in terms of $G,M,d$. (F=ma)
\end{exercise}
\begin{exercise}[medium]
Using $U_{\rm eff}$, explain why no orbit with $L\ne0$ can have
$r\to0$. (concept)
\end{exercise}
